% !PRINTED 548
\weberpage{548}
%<<CONT
I. \gesperrt{Summe und Differenz irgend zweier Zahlen in $\mathfrak{a}$ geben immer wieder Zahlen in $\mathfrak{a}$.}

II. \gesperrt{Das Product irgend einer Zahl in $\mathfrak{a}$ und einer Zahl in $\mathfrak{o}$ gehört dem System $\mathfrak{a}$ an.}

Dieser Forderung würde das aus der einzigen Zahl Null bestehende System genügen, was aber der Einfachheit halber nicht als ein Ideal bezeichnet wird.

Das System $\mathfrak{o}$ dagegen ist ein eigentliches Ideal. Ebenso ist das System aller durch eine bestimmte Zahl $\mu$ in $\mathfrak{o}$ theilbaren Zahlen $\mathfrak{o}\mu$ ein Ideal, und ein solches wird ein \gesperrt{Hauptideal} genannt.

Unter dem Producte $\mathfrak{a}\mathfrak{b}$ zweier Ideale $\mathfrak{a}$ und $\mathfrak{b}$ versteht man den Inbegriff aller Zahlen, die man erhält, wenn man irgend eine Zahl $\alpha$ aus $\mathfrak{a}$ mit einer Zahl $\beta$ aus $\mathfrak{b}$ multiplicirt und eine beliebige Anzahl solcher Zahlenproducte addirt, also den Inbegriff aller Zahlen von der Form $\Sigma \alpha\beta$. Dass dieses Product $\mathfrak{a}\mathfrak{b}$ wieder ein Ideal ist, leuchtet unmittelbar ein. Nach dieser Definition ist z. B. $\mathfrak{o}\mathfrak{a} = \mathfrak{a}$, und das Ideal $\mathfrak{o}$ spielt bei dieser Multiplication die Rolle der Einheit.

Man kann nun die Ideale und Functionale in der Weise auf einander beziehen, dass dabei folgende Gesetze obwalten:

1. \gesperrt{Jedem ganzen Functional entspricht ein bestimmtes Ideal, und associirten Functionalen entspricht dasselbe Ideal.}

2. \gesperrt{Jedem Ideal entsprechen unendlich viele, aber nur associirte ganze Functionale.}

3. \gesperrt{Dem Producte zweier oder mehrerer ganzer Functionale entsprechen die Producte der den Factoren entsprechenden Ideale.}

4. \gesperrt{Einer ganzen Zahl entspricht ein Hauptideal.}

5. \gesperrt{Den Einheiten entspricht das Ideal $\mathfrak{o}$.}

Um dieses Entsprechen zu definiren, ordnen wir zunächst dem Systeme aller Einheiten das Ideal $\mathfrak{o}$ zu. Ist dann ferner $\varphi$ irgend ein ganzes Functional, was keine Einheit ist, so genügt der Inbegriff aller durch $\varphi$ theilbaren ganzen Zahlen $\alpha$ des Körpers $\Omega$ nach \S.\,138 den Forderungen I, II, und ist also ein Ideal, das wir mit $\mathfrak{a}$ bezeichnen\footnote{Zur Bezeichnung der Ideale gebrauchen wir mit \gesperrt{Dedekind} die kleinen \gesperrt{deutschen Buchstaben}.} und dem Functional $\varphi$ zu-
%CONT>>
